\section{Related Work\label{sec:work}}

Architecture recovery and code-to-architecture mapping both infer a system's modular structure from implementation artifacts, such as static dependencies, code semantics, and directory organization. Recovery produces a decomposition without a reference architecture, while mapping assigns entities to the modules of a given one. We evaluate mapping and draw on recovery research for representations of dependencies, lexical content, and project organization.

Recovery techniques have historically been built around one dominant signal: dependency graphs, partitioned by search-based optimization~\cite{bunch,bunch_sahc_nahc,dagc,wca,mca_eca} or clustering~\cite{cct,hcc_2019,limbo,fca,neighborhood_tree_2019,gma_2021}, identifiers and comments~\cite{arc,zbr,risi2012fold,kargar2017semantic}, or file names and paths~\cite{file_clustering}. Each signal is noisy in isolation, since shared utilities introduce cross-module edges and reused terminology mixes entities from many modules, so hybrid methods combine sources~\cite{sc_misra,escrc,nsga3,sadat2019multi,ensemble_recovery,evol,sade,sarif}.

Mapping began with orphan adoption~\cite{tzerpos1997orphan} and the attraction functions of HuGMe~\cite{christl2005equipping,christl2007automated}, and later varied the signal and the learner: textual similarity~\cite{bittencourt_IR}, supervised classifiers~\cite{bibi2016supervised}, and semi-supervised Naive Bayes combining lexical and dependency-derived information~\cite{olsson2019semi,olsson2022automatically}. Text-based approaches compare classifiers for module assignment~\cite{florean2021textclassifierscomparison}, recommend mappings from natural-language architecture descriptions~\cite{InMap}, or integrate package- and class-level granularity~\cite{sinkala2023integrated}. More recent work uses TF-IDF, GloVe, and LLM representations~\cite{johansson2024mapping}, extracts architectural concepts from documentation and code~\cite{fuchss2025enabling}, or prompts LLMs to map implementation units to a reference architecture, which faces context-size limits and repeated processing costs on large codebases~\cite{rukmono_deductive}. Johansson and Caporuscio map functions and files to textual module descriptions with pre-trained embeddings and find files classified more reliably than functions, with performance associated with vocabulary overlap and class separability~\cite{johansson2025mapping}.

Representation-learning work on recovery encoded one or two sources at a time. NEGAR clusters Node2Vec embeddings of the dependency graph alone~\cite{negar}. Jabeen et al.~\cite{structsar} study dependencies alone and with folder organization, comparing random walk-based methods and variational graph autoencoders. GAER~\cite{gaer}, a graph autoencoder, and ARL~\cite{arl_2026}, a contrastively trained GAT, each fix folder and identifier tokens into one encoding without comparing alternatives. For mapping, GNNs were evaluated on homogeneous and dependency-type-aware graphs without folder or member nodes and without separating graph contribution from feature contribution~\cite{jabeen2025graph}. Its extension combines Word2Vec identifier embeddings and binary bag-of-words representations of folder path tokens with dependencies in one model~\cite{c2a_mapping}.

These approaches motivate our distinction between representing an information source and using it during prediction. Text-based mapping reports weaker results for functions than for files~\cite{johansson2026mapping}, but does not examine whether member information helps as context for the file. Our comparisons span lexical, structural, and contextual inputs but retain differences in their training procedures and feature content, so the matched dependency controls and directory-error analyses assess specific design choices and do not turn the broader comparison into an isolated ranking of information sources.